% Subsection 2.3.3: CoT 在代码生成中的应用

算法代码生成具有严格语法约束、强逻辑依赖、需精确处理边界条件等特点，属于典型的高度结构化生成任务。直接生成往往使模型跳过问题抽象与步骤规划，导致解法整体失效。链式思维通过以下方式改善生成质量：
\begin{itemize}
    \item 显式化推理过程，将“思路设计”与“代码实现”解耦；
    \item 增强模型对复杂逻辑、循环结构、状态维护的处理能力；
    \item 提升边界用例、特殊条件的覆盖率，减少运行时错误；
    \item 与参数高效微调形成互补，实现“参数适配+推理引导”的协同增强。
\end{itemize}

图~\ref{fig:ch2-cot-comparison} 对比了直接生成（Direct）、零样本链式思维（CoT-ZS）与少样本链式思维（CoT-FS）三种推理结构在 prompt 组织上的核心差异：Direct 仅向模型传入题目本体并约束「只输出代码」；CoT-ZS 通过 system 指令显式触发分步推理后再生成代码；CoT-FS 在此基础上前置 $k$ 条「题目—推理—代码」示例，借助 in-context learning 进一步约束推理形态。三者在第四章对应推理侧三种 prompt 结构，与训练侧 LoRA 共同构成析因实验的两条独立维度。

\begin{figure}[!htb]
  \centering
  \begin{tikzpicture}[
    node distance=3mm and 18mm,
    col/.style={draw, rounded corners=2pt, align=center, text width=36mm, font=\scriptsize, inner sep=3pt},
    sysn/.style={col, fill=black!8},
    usrn/.style={col, fill=blue!8},
    asst/.style={col, fill=green!10},
    ex/.style={col, fill=orange!12, font=\tiny},
    arr/.style={-{Stealth[length=2mm]}, thick},
    title/.style={font=\footnotesize\bfseries, align=center, text width=40mm},
  ]
    % --- Direct ---
    \node[title] (tD) {Direct\\(单阶段)};
    \node[sysn, below=of tD] (sD) {System：\\You are an expert programmer.};
    \node[usrn, below=of sD] (uD) {User：题目本体\\+「only output the code」};
    \node[asst, below=of uD] (aD) {Assistant：\\\texttt{```python ... ```}};
    \draw[arr] (sD) -- (uD);
    \draw[arr] (uD) -- (aD);

    % --- CoT-ZS ---
    \node[title, right=of tD] (tZ) {CoT-ZS\\(两阶段)};
    \node[sysn, below=of tZ] (sZ) {System：expert + 要求先编号推理后输出 Python};
    \node[usrn, below=of sZ] (uZ) {User：题目本体\\+「Think step by step」};
    \node[asst, below=of uZ] (rZ) {Stage-1 输出：\\编号推理步骤};
    \node[asst, below=of rZ] (aZ) {Stage-2 输出：\\\texttt{```python ... ```}};
    \draw[arr] (sZ) -- (uZ);
    \draw[arr] (uZ) -- (rZ);
    \draw[arr] (rZ) -- node[midway, right, font=\tiny, xshift=1mm]{拼接推理后续写} (aZ);

    % --- CoT-FS ---
    \node[title, right=of tZ] (tF) {CoT-FS$(k)$\\(两阶段 + $k$ 条示例)};
    \node[sysn, below=of tF] (sF) {System：同 CoT-ZS};
    \node[ex, below=of sF] (e1) {Few-shot 1：\\User 题 / Assistant 推理+代码};
    \node[ex, below=of e1] (eN) {$\cdots$ Few-shot $k$};
    \node[usrn, below=of eN] (uF) {User：目标题\\(模板同 CoT-ZS)};
    \node[asst, below=of uF] (rF) {Stage-1：推理};
    \node[asst, below=of rF] (aF) {Stage-2：代码};
    \draw[arr] (sF) -- (e1);
    \draw[arr] (e1) -- (eN);
    \draw[arr] (eN) -- (uF);
    \draw[arr] (uF) -- (rF);
    \draw[arr] (rF) -- (aF);
  \end{tikzpicture}
  \caption{三种推理侧 prompt 结构对比：Direct 单阶段直出代码；CoT-ZS 通过 system 指令触发两阶段推理（先输出推理、再拼接续写代码）；CoT-FS 在 CoT-ZS 基础上前置 $k$ 条「题目—推理—代码」对话示例}
  \label{fig:ch2-cot-comparison}
\end{figure}

在面向算法题的代码生成任务中，CoT 不仅能提升可执行正确率，还能让模型行为更可控、更易解释，是提升复杂场景稳定性的关键技术，也构成本文的核心研究内容之一。